% ===== PRÉAMBULE CANONIQUE — à coller VERBATIM en tête de CHAQUE .tex =====
\documentclass[11pt,a4paper]{article}
\usepackage[utf8]{inputenc}\usepackage[T1]{fontenc}\usepackage{lmodern}
\usepackage[french]{babel}
\usepackage{amsmath,amssymb,mathtools}\usepackage{icomma}
\usepackage[version=4]{mhchem}          % \ce{...} : équations & formules
\usepackage{siunitx}\sisetup{locale=FR} % \qty{}{}, \si{}, \ang{}
\usepackage{chemfig}                    % \chemfig{...} : structures développées
\usepackage{xcolor}\usepackage{colortbl}
\usepackage{tikz}\usepackage{pgfplots}\pgfplotsset{compat=1.18}
\usepackage[most]{tcolorbox}\usepackage{enumitem}\usepackage{array,booktabs}
\usepackage[a4paper,margin=2.2cm]{geometry}
\usetikzlibrary{arrows.meta,calc,angles,quotes,positioning,patterns,backgrounds,shapes.geometric}
\definecolor{primary}{RGB}{222,49,99}\definecolor{primarydark}{RGB}{150,22,55}
\definecolor{defcol}{RGB}{0,114,178}\definecolor{methcol}{RGB}{52,92,148}
\definecolor{excol}{RGB}{0,135,90}\definecolor{attncol}{RGB}{200,40,55}
\tcbset{boxbase/.style={enhanced,breakable,boxrule=0.9pt,arc=2.5pt,
  left=3mm,right=3mm,top=2.3mm,bottom=2.3mm,fonttitle=\bfseries,coltitle=white}}
% --- boîtes TUTORIEL (noms EXACTS, ne pas renommer) ---
\newtcolorbox{cle}[1][]{boxbase,colback=defcol!6,colframe=defcol,
  title={\textsc{À retenir}\ifx\relax#1\relax\relax\else\textmd{ : \itshape #1}\fi}}
\newtcolorbox{methode}[1][]{boxbase,colback=methcol!7,colframe=methcol,sharp corners,
  title={\textsc{Méthode}\ifx\relax#1\relax\relax\else\textmd{ --- \itshape #1}\fi}}
\newtcolorbox{exemplebox}[1][]{boxbase,colback=excol!5,colframe=excol,
  title={\textsc{Exemple}\ifx\relax#1\relax\relax\else\textmd{ : \itshape #1}\fi}}
\newtcolorbox{piege}[1][]{boxbase,colback=attncol!6,colframe=attncol,
  title={\textsc{Piège}\ifx\relax#1\relax\relax\else\textmd{ : \itshape #1}\fi}}
\newtcolorbox{remarque}[1][]{boxbase,colback=black!4,colframe=black!45,
  title={\textsc{Remarque}\ifx\relax#1\relax\relax\else\textmd{ : \itshape #1}\fi}}
% --- boîtes EXERCICE / CORRIGÉ (noms EXACTS, auto-comptées) ---
\newtcolorbox[auto counter]{exo}[1][]{boxbase,colback=primary!4,colframe=primary,
  title={\textsc{Exercice}~\thetcbcounter\ifx\relax#1\relax\relax\else\textmd{ --- \itshape #1}\fi}}
\newtcolorbox[auto counter]{corrige}[1][]{boxbase,colback=excol!4,colframe=excol,
  title={\textsc{Corrigé}~\thetcbcounter\ifx\relax#1\relax\relax\else\textmd{ --- \itshape #1}\fi}}
\newcommand{\R}{\mathbb{R}}\newcommand{\N}{\mathbb{N}}\newcommand{\Z}{\mathbb{Z}}\newcommand{\ds}{\displaystyle}
% =========================================================================
\begin{document}
\sisetup{per-mode=symbol}
\begin{corrige}[quel phosphate libère le plus d'ions phosphate ?]
\begin{enumerate}[label=\alph*)]
  \item Oui : tous de type $3{:}2$ ($\ce{M3(PO4)2}$), donc $K_{\mathrm{ps}}=108\,s^{5}$ pour chacun.
        Comme ils ont \emph{la même} relation, on peut cette fois comparer directement leurs
        $K_{\mathrm{ps}}$.
  \item $[\ce{PO4^{3-}}]=2s$, avec $s=\sqrt[5]{K_{\mathrm{ps}}/108}$. Donc $[\ce{PO4^{3-}}]$ est une
        fonction \textbf{croissante} de $K_{\mathrm{ps}}$.
  \item Le plus grand $K_{\mathrm{ps}}$ donne le plus grand $s$, donc le plus de \ce{PO4^{3-}} : c'est
        \textbf{\ce{Sr3(PO4)2}} ($K_{\mathrm{ps}}=\num{1.0e-31}$, le plus élevé).
\end{enumerate}
\end{corrige}

\begin{corrige}[le piège de la comparaison inter-types]
\begin{enumerate}[label=\alph*)]
  \item \textbf{Non} : les trois sels sont de types \emph{différents}, donc on ne peut pas comparer
        leurs $K_{\mathrm{ps}}$ directement. Il faut repasser par $s$.
  \item Calculs :
    \begin{itemize}[leftmargin=1.6em,topsep=1pt,itemsep=1pt]
      \item \ce{AgCl} : $s=\sqrt{\num{1.8e-10}}=\qty{1.3e-5}{\mole\per\litre}$ ;
      \item \ce{Ag2CrO4} : $s=\sqrt[3]{\num{1.1e-12}/4}=\sqrt[3]{\num{2.75e-13}}=\qty{6.5e-5}{\mole\per\litre}$ ;
      \item \ce{CaF2} : $s=\sqrt[3]{\num{3.9e-11}/4}=\sqrt[3]{\num{9.75e-12}}=\qty{2.14e-4}{\mole\per\litre}$.
    \end{itemize}
  \item Par solubilité croissante :
    $\ce{AgCl}\ (\num{1.3e-5}) < \ce{Ag2CrO4}\ (\num{6.5e-5}) < \ce{CaF2}\ (\num{2.14e-4})$.
    \textbf{Le piège :} \ce{Ag2CrO4}, qui a le plus \emph{petit} $K_{\mathrm{ps}}$, n'est pas le moins
    soluble — c'est \ce{AgCl}. Comparer des $K_{\mathrm{ps}}$ de types différents est trompeur.
\end{enumerate}
\end{corrige}

\begin{corrige}[$K_{\mathrm{ps}}$ à partir de la solubilité massique (forme littérale)]
\begin{enumerate}[label=\alph*)]
  \item Type $1{:}2$ : $K_{\mathrm{ps}}=4s^{3}$ avec $s=s_m/M$, d'où
    \[ K_{\mathrm{ps}} = 4\left(\dfrac{s_m}{M}\right)^{3} = 4\left(\dfrac{\qty{4.84}{\gram\per\litre}}{M}\right)^{3}. \]
  \item Avec $M=\qty{367.0}{\gram\per\mole}$ : $s=\qty{1.32e-2}{\mole\per\litre}$ et
    \[ K_{\mathrm{ps}} = 4\,(\num{1.32e-2})^{3} = \num{9.2e-6}. \]
\end{enumerate}
\end{corrige}

\begin{corrige}[de bout en bout : le fluorure de calcium]
\begin{enumerate}[label=\alph*)]
  \item Type $1{:}2$ : $s=\sqrt[3]{\dfrac{K_{\mathrm{ps}}}{4}}=\sqrt[3]{\dfrac{\num{3.9e-11}}{4}}
        =\sqrt[3]{\num{9.75e-12}}=\qty{2.14e-4}{\mole\per\litre}$.
  \item $[\ce{Ca^{2+}}]=s=\qty{2.14e-4}{\mole\per\litre}$ et
        $[\ce{F-}]=2s=\qty{4.28e-4}{\mole\per\litre}$.
  \item $s_m=s\times M=\qty{2.14e-4}{\mole\per\litre}\times\qty{78.1}{\gram\per\mole}
        =\qty{1.67e-2}{\gram\per\litre}$.
\end{enumerate}
\end{corrige}
\end{document}
